\documentclass[12pt,a4paper]{article}
\usepackage{amsmath,amssymb,mathrsfs,tikz,times,pifont}
\usepackage{enumitem}
\newcommand\circitem[1]{%
\tikz[baseline=(char.base)]{
\node[circle,draw=gray, fill=red!55,
minimum size=1.2em,inner sep=0] (char) {#1};}}
\newcommand\boxitem[1]{%
\tikz[baseline=(char.base)]{
\node[fill=cyan,
minimum size=1.2em,inner sep=0] (char) {#1};}}
\setlist[enumerate,1]{label=\protect\circitem{\arabic*}}
\setlist[enumerate,2]{label=\protect\boxitem{\alph*}}
%%%::::::by chnini ameur :::::::%%%
\everymath{\displaystyle}
\usepackage[left=1cm,right=1cm,top=1cm,bottom=1.7cm]{geometry}
\usepackage{array,multirow}
\usepackage[most]{tcolorbox}
\usepackage{varwidth}
\tcbuselibrary{skins,hooks}
\usetikzlibrary{patterns}
%%%::::::by chnini ameur :::::::%%%
\newtcolorbox{exa}[2][]{enhanced,breakable,before skip=2mm,after skip=5mm,
colback=yellow!20!white,colframe=black!20!blue,boxrule=0.5mm,
attach boxed title to top left ={xshift=0.6cm,yshift*=1mm-\tcboxedtitleheight},
fonttitle=\bfseries,
title={#2},#1,
% varwidth boxed title*=-3cm,
boxed title style={frame code={
\path[fill=tcbcolback!30!black]
([yshift=-1mm,xshift=-1mm]frame.north west)
arc[start angle=0,end angle=180,radius=1mm]
([yshift=-1mm,xshift=1mm]frame.north east)
arc[start angle=180,end angle=0,radius=1mm];
\path[left color=tcbcolback!60!black,right color = tcbcolback!60!black,
middle color = tcbcolback!80!black]
([xshift=-2mm]frame.north west) -- ([xshift=2mm]frame.north east)
[rounded corners=1mm]-- ([xshift=1mm,yshift=-1mm]frame.north east)
-- (frame.south east) -- (frame.south west)
-- ([xshift=-1mm,yshift=-1mm]frame.north west)
[sharp corners]-- cycle;
},interior engine=empty,
},interior style={top color=yellow!5}}
%%%%%%%%%%%%%%%%%%%%%%%

\usepackage{fancyhdr}
\usepackage{eso-pic}         % Pour ajouter des éléments en arrière-plan
% Commande pour ajouter du texte en arrière-plan
\AddToShipoutPicture{
    \AtTextCenter{%
        \makebox[0pt]{\rotatebox{80}{\textcolor[gray]{0.7}{\fontsize{5cm}{5cm}\selectfont PGB}}}
    }
}
\usepackage{lastpage}
\fancyhf{}
\pagestyle{fancy}
\renewcommand{\footrulewidth}{1pt}
\renewcommand{\headrulewidth}{0pt}
\renewcommand{\footruleskip}{10pt}
\fancyfoot[R]{
\color{blue}\ding{45}\ \textbf{2025}
}
\fancyfoot[L]{
\color{blue}\ding{45}\ \textbf{Prof:M. BA}
}
\cfoot{\bf
\thepage /
\pageref{LastPage}}
\begin{document}
\renewcommand{\arraystretch}{1.5}
\renewcommand{\arrayrulewidth}{1.2pt}
\begin{tikzpicture}[overlay,remember picture]
\node[draw=blue,line width=1.2pt,fill=purple,text=blue,inner sep=3mm,rounded corners,pattern=dots]at ([yshift=-2.5cm]current page.north) {\begingroup\setlength{\fboxsep}{0pt}\colorbox{white}{\begin{tabular}{|*1{>{\centering \arraybackslash}p{0.28\textwidth}} |*2{>{\centering \arraybackslash}p{0.2\textwidth}|} *1{>{\centering \arraybackslash}p{0.19\textwidth}|} }
\hline
\multicolumn{3}{|c|}{$\diamond$$\diamond$$\diamond$\ \textbf{Lycée de Dindéfélo}\ $\diamond$$\diamond$$\diamond$ }& \textbf{A.S. : 2023/2024} \\ \hline
\textbf{Matière: Mathématiques}& \textbf{Niveau : T}\textbf{S2} &\textbf{Date: 08/05/2024} & \textbf{Durée : 4 heures} \\ \hline
\multicolumn{4}{|c|}{\parbox[c]{10cm}{\begin{center}
\textbf{{\Large\sffamily Devoir n$ ^{\circ} $ 1 Du 1$ ^\text{\bf er} $ Semestre}}
\end{center}}} \\ \hline
\end{tabular}}\endgroup};
\end{tikzpicture}
\vspace{3cm}

% Exercice 1
\subsection*{Exercice 1 \mdseries(1,25 point)}
Calculer les limites suivantes : \textbf{(0,75 pt + 0,5 pt)}
\[
\lim_{x \to +\infty} \frac{\ln(x-1)}{2x-3} \quad\quad \lim_{x \to +\infty} \frac{\ln(x)}{\sqrt{x}}
\]

% Exercice 2
\subsection*{Exercice 2 \mdseries(2,5 points)}
\begin{enumerate}
  \item Résoudre dans $\mathbb{C}$ l'équation $z^{4}=1$. \textbf{(1 pt)}
  \item En déduire les solutions dans $\mathbb{C}$ de l'équation $\left[\left(\sqrt{3}-i\right)z+i\right]^{4}=1$. \textbf{(1,5 pt)}
\end{enumerate}

% Exercice 3
\subsection*{Exercice 3 \mdseries(4 points)}
Le plan complexe est muni d’un repère orthonormé $(O, \vec{u}, \vec{v})$.

On considère les points $A$ et $B$ d’affixes respectives $a$ et $b$ telles que :
\[
a = -\left(\frac{1+\sqrt{3}}{2}\right)(1-i) \quad \text{et} \quad b = -\left(\frac{1+\sqrt{3}}{2}\right)(1+i)
\]
\begin{enumerate}
  \item Écrire $a$ et $b$ sous forme trigonométrique. \textbf{(1 pt)}
  \item Quelle est la nature du triangle $OAB$ ? \textbf{(1 pt)}
  \item On désigne par $C$ le point d’affixe $1$.
  \begin{enumerate}
    \item Écrire le nombre complexe $\dfrac{1-a}{b-a}$ sous forme trigonométrique. \textbf{(1 pt)}
    \item En déduire la nature du triangle $ABC$ puis utiliser ce résultat pour placer les points $A$, $B$, $C$ sur la figure. \textbf{(1 pt)}
  \end{enumerate}
\end{enumerate}

% Problème
\subsection*{Problème \mdseries(12,25 points)}

\subsubsection*{Partie A \mdseries(5,5 points)}
Soit $f$ la fonction numérique définie par $f(x) = -x + e^{x-1}$ et $\mathcal{C}_{f}$ sa courbe représentative dans un repère orthonormé $(O, \vec{i}, \vec{j})$ (unité graphique : 1 cm).

\begin{enumerate}
  \item Calculer les limites de $f$ aux bornes de son ensemble de définition puis étudier la nature des branches infinies. \textbf{(2,25 pts = 0,25 pt + 0,5 pt $\times$ 4)}
  \item Préciser la position de la courbe $\mathcal{C}_{f}$ par rapport à son asymptote. \textbf{(0,5 pt)}
  \item Dresser le tableau de variation de $f$ puis en déduire le signe de $f(x)$. \textbf{(1 pt = 0,5 pt $\times$ 2)}
  \item Construire $\mathcal{C}_{f}$. \textbf{(0,75 pt)}
  \item Soit $g$ la fonction définie sur $\mathbb{R}$ par $g(x) = x + e^{-x-1}$.
  \begin{enumerate}
    \item Exprimer $g(x)$ en fonction de $f(x)$. \textbf{(0,5 pt)}
    \item Construire alors $\mathcal{C}_{g}$ dans le même repère. \textbf{(0,5 pt)}
  \end{enumerate}
\end{enumerate}

\subsubsection*{Partie B \mdseries(6,75 points)}
Soit $h$ la fonction définie par $h(x) = -x + \ln(-x + e^{x-1})$ et $\Gamma$ sa courbe représentative dans un autre repère orthonormal.

\begin{enumerate}
  \item 
  \begin{enumerate}
    \item Déterminer l’ensemble de définition de $h$ puis calculer les limites de $h$ aux bornes de son ensemble de définition. \textbf{(2,25 pts = 0,25 pt + 0,5 pt $\times$ 4)}\\
    En déduire les asymptotes éventuelles. \textbf{(1 pt = 0,5 pt $\times$ 2)}\\
    Étudier le comportement de $\Gamma$ au voisinage de $-\infty$. \textbf{(0,75 pt)}
    \item Montrer que pour tout $x > 1$, $h(x) = \ln(-x e^{-x} + e^{-1})$. \textbf{(0,25 pt)}
  \end{enumerate}
  \item Dresser le tableau de variation de $h$. \textbf{(0,75 pt)}
  \item Soit $\varphi$ la restriction de $h$ à l’intervalle $I = \left]-\infty, 1\right[$.
  \begin{enumerate}
    \item Montrer que $\varphi$ réalise une bijection de $I$ vers un intervalle $J$ à préciser. \textbf{(0,25 pt)}
    \item Montrer que l’équation $\varphi(x) = 0$ admet une solution unique $\alpha$.\\
    En déduire que $\alpha \in \left]-0{,}5 ; -0{,}4\right[$ puis donner un encadrement à $10^{-2}$ près de $\alpha$. \textbf{(0,25 pt $\times$ 2)}
    \item Construire $\Gamma$ et la courbe de $\varphi^{-1}$ dans le même repère. \textbf{(1 pt = 0,5 pt $\times$ 2)}
  \end{enumerate}
\end{enumerate}

\end{document}